\documentclass[12 pt, letterpaper]{hmcpset}
\usepackage{hw-preamble}

\name{Jayden Thadani}
\class{MATH 145 --- Topics in Geometry and Topology}
\assignment{Homework 11}
\duedate{23rd April 2025}

\begin{document}

\begin{problem}[59.]
	Recall that the $n$-orthoplex is defined to be the simplicial complex with $2n$ vertices
	\[
		V = \{ \pm 1, \dots, \pm n \}
	\]
	whose simplices are the subsets that do not contain any pair $\pm k$. It is a model for the sphere $S^{n - 1}$.

	\begin{enumerate}
		\item List the vertices, edges, and triangles of the $3$-orthoplex.

		\item Write down the matrices for the boundary maps $\partial_{0}, \partial_{1}, \partial_{2}$ and determine their ranks.

		\item Obtain the Betti sequence for the $3$-orthoplex.
	\end{enumerate}
\end{problem}

\begin{solution}
	Denote the $3$-orthoplex by $X$.
	\begin{enumerate}
		\item
			\[
				\boxed{
					\begin{split}
						X^{(0)} & = \{ \{ \pm 1 \}, \{ \pm 2 \}, \{ \pm 3 \} \} \\
						X^{(1)} \setminus X^{(0)} & = \{ \{ \pm 1, \pm 2 \}, \{ \pm 2, \pm 3 \}, \{ \pm 3, \pm 1 \} \} \\
						X^{(2)} \setminus X^{(1)} & = \{ \{ 1, 2, 3 \}, \{ -1, -2, -3 \} \}
					\end{split}
				}
			\]

			We can write the $k$-simplices as $X^{(k)} \setminus X^{(k - 1)}$. We get the sets above by considering all subsets of size $k - 1$ that contain no pairs $\pm m$.

		\item

		\item
			\[
				\boxed{
					\mathbf{b} = (1, 0, 1, 0, 0, \dots)
				}
			\]
	\end{enumerate}
\end{solution}

\pagebreak

\begin{problem}[60.]
	Recall that $\partial B^{n}$ is defined to be the simplicial complex with $n + 1$ vertices
	\[
		V = \{ 0, 1, \dots, n \}
	\]
	whose simplices are the subsets of cardinality at most $n$. It is a model for the sphere $S^{n - 1}$.

	\begin{enumerate}
		\item List the vertices, edges, and triangles of $\partial B^{4}$.

		\item Write down the matrices for the boundary maps $\partial_{0}, \partial_{1}, \partial_{2}$ and determine their ranks.

		\item Obtain the Betti sequence for $\partial B^{4}$.
	\end{enumerate}
\end{problem}

\begin{solution}
	Let $X = \partial B^{4}$
	\begin{enumerate}
		\item
			\[
				\boxed{
					\begin{split}
						X^{(0)} & = \{ \{ 0 \},  \{ 1 \}, \{ 2 \}, \{ 3 \} \} \\
						X^{(1)} \setminus X^{(0)} & = \{ \{ 0, 1 \}, \{ 0, 2 \}, \{ 0, 3 \} \{ 1, 2 \}, \{ 2, 3 \}, \{ 3, 1 \} \} \\
						X^{(2)} \setminus X^{(1)} & = \{ \{ 0, 1, 2 \}, \{ 1, 2, 3 \}, \{ 0, 2, 3 \},  \{ 0, 1, 3 \} \}
					\end{split}
				}
			\]

			Again, we write the set of all $k$-simplices as the difference of skeletons $X^{(k)} \setminus X^{(k - 1)}$.

		\item

		\item
			\[
				\boxed{
					\mathbf{b} = (1, 0, 1, 0, 0, \dots)
				}
			\]
	\end{enumerate}
\end{solution}

\pagebreak

\begin{problem}[61. (the Morse inequalities)]
	Prove the following inequality for a finite simplicial complex:
	\[
		b_{k} - b_{k - 1} + \dots + (-1)^{k} b_{0} \le c_{k} - c_{k - 1} + \dots + (-1)^{k} c_{0}.
	\]
\end{problem}

\begin{solution}
	Recall that for the simplicial chain complex
	\[
	\begin{tikzcd}
		0 & {C}_{0}(X, \mathbb{F}) \ar[l] & {C}_{1}(X, \mathbb{F}) \ar[l, "\partial_{0}"] & {C}_{2}(X, \mathbb{F}) \ar[l, "\partial_{1}"] & \dots \ar[l, "\partial_{2}"]
	\end{tikzcd}
	\]
	we have
	\[
		b_{i} = \dim (\ker \partial_{i - 1} / \operatorname{img} \partial_{i}) = \operatorname{null} \partial_{i - 1} - \operatorname{rank} \partial_{i}
	\]
	by definition, and
	\[
		c_{i} = \operatorname{rank} \partial_{i - 1} + \operatorname{null} \partial_{i - 1}
	\]
	by the rank-nullity theorem.

	Note that
	\begin{align*}
		b_{k} - b_{k - 1} & = \operatorname{null} \partial_{k - 1} - \operatorname{rank} \partial_{k} + \operatorname{null} \partial_{k - 2}
	\end{align*}
	and
	\begin{align*}
		b_{i} - \operatorname{null} \partial_{i} & = \operatorname{null} \partial_{i - 1} - \operatorname{rank} \partial_{i} - \operatorname{null} \partial_{i} \\
							 & = \operatorname{null} \partial_{i - 1} - c_{i}.
	\end{align*}
	We can use this to turn the sum into something like a telescoping sum. We have
	\begin{align*}
		\sum_{i = 0}^{k} (-1)^{k - i} b_{i} & = - \operatorname{rank} \partial_{k} + c_{k}  - \operatorname{null} \partial_{k - 2} + \sum_{i = 0}^{k - 2} (-1)^{k - i} b_{i} \\
						& = - \operatorname{rank} \partial_{k} + c_{k} + b_{k - 2} - \operatorname{null} \partial_{k - 2} + \sum_{i = 0}^{k - 3} (-1)^{k - i} b_{i} \\
						& = -\operatorname{rank} \partial_{k} + c_{k} - c_{k - 1} + \operatorname{null} \partial_{k - 3} + \sum_{i = 0}^{k - 3} \\
						& = \vdots \\
						& = -\operatorname{rank} \partial_{k} + \sum_{i = 0}^{k} (-1)^{k - i} c_{k}.
	\end{align*}

	That is,
	\[
		\sum_{i = 0}^{k} (-1)^{k - i} b_{i} = - \operatorname{rank} \partial_{k} + \sum_{i = 0}^{k} (-1)^{k - i} c_{i} \le \sum_{i = 0}^{k} (-1)^{k - i} c_{i}.
	\]
	If $X$ is $k$-dimensional, then all $n > k$ have $C_{n}(X, \mathbb{F}) = 0$. In particular $C_{k + 1}(X, \mathbb{F}) = 0$ and thus, $\partial_{k} : C_{k + 1}(X, \mathbb{F}) \to C_{k}(X, \mathbb{F})$ has $\operatorname{rank} \partial_{k} = 0$. This is exactly Euler's formula.
\end{solution}

\pagebreak

\begin{problem}[62.]
	The torus is a topological space which can be constructed from a closed square by gluing each side to its opposite side. Here is a representation of a torus as a simplicial complex with $9$ vertices, $27$ edges, and $18$ triangles.
	\begin{figure}[H]
		\centering
		\includegraphics[width = 0.3 \textwidth]{figures/P62 torus}
	\end{figure}

	By calculating the ranks of $\partial_{0}, \partial_{1}$, determine the Betti sequence for the torus.
\end{problem}

\begin{solution}
	\[
		\boxed{
			\mathbf{b} = (1, 2, 1, 0, 0, \dots)
		}
	\]
\end{solution}

\pagebreak

\begin{problem}[63.]
	Let $X = \{ 2, 3, 4, \dots, 16 \}$. Define a simplicial complex $S_{16}$ as follows: a subset $\sigma \subseteq X$ belongs to $S_{16}$ if and only if the elements of $\sigma$ share a common factor greater than $1$.
	\begin{enumerate}
		\item Explain why $S_{16}$ is a simplicial complex.

		\item How many vertices and edges in $S_{16}$?

		\item What is the largest-dimensional simplex in $S_{16}$?

		\item How many connected components does $S_{16}$ have?

		\item Find three edges $\{ a, b \}, \{ b, c \}, \{ c, a \}$ in $S_{16}$ such that $\{ a, b, c \} \notin S_{16}$.

		\item What are the \emph{maximal} simplices of $S_{16}$.

		\item Sketch $S_{16}$ and write down your best guess for its Betti sequence.

		\item The maximal simplices form a cover of the set $X$. Determine and sketch the nerve of this cover.

		\item Calculate the Betti sequence of the nerve. Compare it with your guess in (g).
	\end{enumerate}
\end{problem}

\begin{solution}
	\begin{enumerate}
		\item Suppose $\sigma \in S_{16}$ with all elements sharing a common factor $d > 1$ --- we have $d \mid x \in \sigma$ for every $x \in \sigma$. Suppose $\tau \subseteq \sigma$. Then $d \mid x$ for every $x \in \tau$. Thus, all elements of $\tau$ share a common factor greater than $1$.

			This suffices to show $S_{16}$ is a simplicial complex.

		\item 
			\[
				\boxed{
					\begin{split}
						S_{16}^{(0)} & = \{ \{ 2 \}, \{ 3 \}, \dots, \{ 16 \} \} \\
						S_{16}^{(1)} \setminus S_{16}^{(0)} & = \binom{X}{2} \setminus \\
										    & \left (\{ \{ n, p \} \mid p \text{ prime} \} \cup \{ \{ 4, 9 \}, \{ 4, 15 \}, \{ 8, 9 \}, \{ 8, 15 \}, \{ 9, 10 \}, \{ 9, 14 \}, \{ 9, 16 \}, \{ 14, 15 \}, \{ 15, 16 \} \} \right ) . \\
						S_{16}^{(2)} \setminus S_{16}^{(1)} & = \{ \{ m_{1}, m_{2}, m_{3} \} \mid m_{1}, m_{2}, m_{3} \in 2 \mathbb{Z} \} \cup \{ \{ m_{1}, m_{2}, m_{3} \} \mid m_{1}, m_{2}, m_{3} \in 3 \mathbb{Z} \} \cup \{ \{ 5, 10, 15 \} \}
					\end{split}
				}
			\]

			All of the singleton sets are vertices --- they are all greater than one, and are thus, divisible by themselves, a number greater than one.

			For the edges, we can ignore all pairs where one number is prime as well as the pairs of co-prime numbers listed above.

			For triangles we break the triples down by the common prime dividing them. There are $\binom{8}{3}$ divided by $2$, $\binom{5}{3}$ divided by $3$, and $1$ divided by $5$.

		\item
			\[
				\boxed{
					\begin{split}
						\sigma & = \{ 2, 4, 6, \dots, 16 \} \\
						\dim \sigma & = 7.
					\end{split}
				}
			\]

			The largest-dimensional simplex of $S_{16}$ consists of the $8$ even numbers less than $16$.

		\item
			\[
				\boxed{
					b_{0}(X) = 3.
				}
			\]

			Every number that is not a prime greater than $8$ is divisible by a number less than $8$. That number divides an even number. Thus, almost all numbers are connected to the largest simplex of all even numbers.

			The primes greater than $8$, namely $11$ and $13$ share common factors with none of the other numbers, and thus, form two more connected components.

		\item 
			\[
				\boxed{
					\{ 6, 10, 15 \} \notin S_{16}, \text{ but } \{ 6, 10 \}, \{ 10, 15 \}, \{ 15, 6 \} \in S_{16}.
				}
			\]

		\item
			\[
				\boxed{
					\mathcal{U} = \{ U_{p} = \{ pk \mid p \text{ prime}, k \in \mathbb{N} \} \}
				}
			\]
			
			Every simplex $\sigma$ of $S_{16}$ is a set of numbers with a common divisor $d$ greater than $1$. Then the smallest prime $p$ dividing $d$ divides all the numbers in the simplex. Thus, the simplex corresponding to all multiples of $p$, denoted $U_{p}$, contains $\sigma$. 

			These $U_{p}$ are maximal since they include the prime $p$ which is only divisible by itself, and is thus, contained in no larger simplex.

		\item
			\[
				\boxed{
					\mathbf{b} = (3, 1, 0, 0, \dots)
				}
			\]

			\begin{figure}[H]
				\centering
				\includegraphics[width = 0.5 \textwidth]{figures/P63g guess.pdf}
				\caption{$S_{16}$ drawn with $7$ simplices highlighted in red, $5$-simplices highlighted in yellow, $2$-simplices highlighted in black, and edges in black, only drawn outside of simplices}
			\end{figure}

			Clearly this has $3$ connected components.

			Notice that after contracting all the large $k$-simplices onto the drawn edges, we are left with a $1$-cycle $[6, 10] + [10, 15] + [15, 6]$ that is not the boundary of any $2$-chain. Thus, we have exactly one hole.
		\item
			\[
				\boxed{
					\operatorname{nerve} \mathcal{U} = \{ \{ U_{2} \}, \{ U_{3} \}, \{ U_{5} \}, \{ U_{7} \}, \{ U_{11} \}, \{ U_{13} \}, \{ U_{2}, U_{3} \}, \{ U_{2}, U_{5} \}, \{ U_{2}, U_{7} \}  \{ U_{3}, U_{5} \} \}.
				}
			\]
			
			\begin{figure}[H]
				\centering
				\includegraphics[width = 0.5 \textwidth]{figures/P63h nerve.pdf}
				\caption{$\operatorname{nerve} \mathcal{U}$}
			\end{figure}

			Since the condition that two maximal simplicial complexes intersect is equivalent to the condition that their least common multiple is at most $16$, we can write the nerve as
			\[
				\operatorname{nerve} \mathcal{U} = \{ \{ U_{p_{1}}, \dots, U_{p_{n}} \} \mid p_{1} \cdots p_{n} \le 16 \}.
			\]
			(The least common multiple of distinct primes is just their product). There are no triples of primes such that this is true (the three smallest primes have product $2 \times 3 \times 5 = 30 > 16$). The pairs of primes for which this is true are $2 \times 3 = 6$, $2 \times 5 = 10$, $2 \times 7 = 14$ and $3 \times 5 = 15$. Thus,
			\[
				\operatorname{nerve} \mathcal{U} = \{ \{ U_{2} \}, \{ U_{3} \}, \{ U_{5} \}, \{ U_{7} \}, \{ U_{11} \}, \{ U_{13} \}, \{ U_{2}, U_{3} \}, \{ U_{2}, U_{5} \}, \{ U_{2}, U_{7} \}  \{ U_{3}, U_{5} \} \}.
			\]

		\item
			\[
				\boxed{
					\mathbf{b} = (3, 1, 0, 0, \dots)
				}
			\]

			We've already shown that $b_{0} = 3$. In the sketch of the nerve there is only one $1$-cycle, and it's not a $1$-boundary. This is the unique non-trivial dimension of the homology vector space. Thus, $b_{1} = 1$.
			The nerve is $1$-dimensional, so has trivial higher homology.

			The Betti number of the nerve is exactly the same as the estimated Betti number of $S_{16}$.
	\end{enumerate}
\end{solution}

\end{document}
